\section{Related Work}\label{sec:related}

\paragraph{Training acceleration for deep learning.}
Reducing wall-clock training time without sacrificing accuracy is a longstanding goal.
The DAWNBench competition~\cite{coleman2017dawnbench} formalised this as a benchmark, spurring work on cyclic learning rates~\cite{smith2019superconvergence}, mixed-precision training, and large-batch {SGD}.
Garg~\etal~\cite{garg2024cifar10} pushed CNN classification to $94\%$ on \cifar{} in 3.29 seconds using four complementary tricks (patch whitening, Dirac init, ScaleBias+Lookahead, alternating flip).
Our work investigates whether these same acceleration techniques remain effective when the model architecture (CNN $\to$ ViT) or task (classification $\to$ reconstruction) changes.

\paragraph{Vision Transformers on small datasets.}
The original ViT~\cite{dosovitskiy2021vit} required large-scale pretraining to outperform CNNs on small datasets.
DeiT~\cite{touvron2021deit} showed that careful training recipes---knowledge distillation, strong augmentation---allow ViTs to match ResNets without pretraining.
Subsequent work~\cite{touvron2022deit3,liu2021swin} confirmed that augmentation strategy is the dominant factor for ViT efficiency on small datasets.
Our dual-recipe design (weak vs.\ strong augmentation) is motivated directly by this finding: we test whether CNN tricks that help under simple augmentation remain effective under the state-of-the-art ViT recipe.

\paragraph{Masked reconstruction.}
Masked image modelling---predicting corrupted or missing pixel regions---has emerged as a powerful self-supervised objective~\cite{he2022masked}.
CNN encoder-decoders were among the first architectures applied to this task.
We study whether CNN initialisation tricks (whitening, Dirac init) that accelerate classification training also accelerate the convergence of a CNN reconstruction model, extending the scope of the transferability question beyond classification.

\paragraph{Initialisation for transformers.}
ReZero~\cite{bachlechner2021rezero} and T-Fixup~\cite{huang2020improve} propose identity-like initialisations for residual transformer blocks, showing that near-identity starts can remove the need for LR warmup.
Our Dirac init is a stronger variant applied to the attention projections themselves; we show that full identity collapses head diversity, while soft blending ($\alpha{=}0.3$) can be mildly beneficial.

\paragraph{Patch whitening.}
Coates~\etal~\cite{coates2011importance} demonstrated that whitened patch features improve shallow CNN representations.
The original ViT uses a learned linear projection; our work is the first to study frozen whitened patch embeddings under strong augmentation, revealing that the frozen-basis/augmentation statistics mismatch---not freezing per se---is the decisive factor.

\paragraph{Lookahead.}
Zhang~\etal~\cite{zhang2019lookahead} introduced Lookahead as a wrapper that stabilises fast-optimizer trajectories via slow-weight interpolation, validated primarily with flat or mildly decaying LR schedules.
We identify the conflict with monotonically decaying schedules (OneCycleLR, cosine decay) as a distinct failure mode independent of architecture.

